\documentclass[10pt,a4paper]{amsart}
\usepackage[margin=19mm]{geometry}
\usepackage{amsmath,amssymb,mathtools}
\usepackage[scr=rsfs,cal=euler]{mathalfa}
\usepackage{xfrac}
\usepackage[unicode,hidelinks,bookmarks=false]{hyperref}
\newcommand{\ie}{i.e.\ }
\newcommand{\cf}{cf.\ }
\newcommand{\resp}{respectively\ }
\newcommand{\Subsection}{\subsection}
\providecommand{\nobreakdash}{\nobreak\hbox{-}}
\setlength{\emergencystretch}{2em}
\multlinegap0pt
\usepackage{amssymb}
\usepackage{tikz}
\usepackage{enumerate}
\newtheorem{lemma}{Lemma}
\title{Covering Complete Geometric Graphs by Monotone Paths}
\author{Adrian Dumitrescu, J\'anos Pach, Morteza Saghafian, Alex Scott}
\date{Source: arXiv:2507.10840v2 (CC BY 4.0)}
\begin{document}
\maketitle

\noindent\emph{Source and licence.} Covering Complete Geometric Graphs by Monotone Paths. Version arXiv:2507.10840v2 (CC BY 4.0), distributed by arXiv under the Creative Commons Attribution 4.0 International licence, http://creativecommons.org/licenses/by/4.0/, as recorded in the arXiv licence metadata for this exact version and shown on the abstract page https://arxiv.org/abs/2507.10840v2. Full licence notice: \texttt{licenses/arxiv-2507.10840-CC-BY-4.0.txt}.\par\smallskip

\noindent\textit{Editorial scope.} Counted unit: the article's lemma lem:locus (source0.tex lines 336--338). The article's complete original proof of it is delivered with it (source0.tex lines 340--350, one \texttt{proof} environment of 11 lines). No mathematics is written, repaired or re-derived: the statement and the proof are the article's own lines, in the article's own notation, and no retained line is substituted or re-flowed.  No further source-local context is needed: every cross-reference of every retained line resolves inside the two delivered spans.  Nothing was omitted from any retained span, so the package carries no omission record.  No retained line contains a citation, so the wrapper carries no bibliography and no bibliography entry.  No retained line contains a figure environment, an image-inclusion command or drawing code, so no image is delivered and the package declares no asset dependency.  Editorial formatting only, in addition: an AMS \texttt{amsart} preamble that restates the article's own declarations and macros used by the retained lines, namely the environments \texttt{lemma} on separate counters, together with the standard packages those lines need.  The retained environments print in this document as Lemma 1 in order of appearance, whereas the article prints the same items with the numbering of its own full document.  The complete native source of the article as distributed in the arXiv e-print for this exact version is bundled unchanged as \texttt{sources/181586/source0.tex}.
\begin{lemma}  \label{lem:locus}
The area of $R(a,b)$ is at most $4 \theta$. 
\end{lemma}

\begin{proof}
  Consider four lines incident to $a$ and $b$ respectively, and making 
(clockwise or counterclockwise) angles of $\theta$ with $ab$.
  The boundary of the locus is made by these four lines and two polygonal arcs
  on the boundary of the square. The locus is easily seen to be contained in an axis-parallel rectangle
  of width $2 \sqrt{2} \tan \theta$ and height $\sqrt2$.
  As such, its area is bounded from above by $ 4 \tan \theta$.
  Excluding the two isosceles triangles based on the left and right vertical sides yields an improved
  area bound of $\frac34 \cdot 4 \tan \theta = 3 \tan \theta \leq 4 \theta$
  (here we assume that $n$ is sufficiently large). 
\end{proof}

\end{document}
